\section{Abstract}\label{sec:3-abstract}
The growing maturity of pulsar timing array experiments — NANOGrav~\cite{afzal_nanograv_2023}, EPTA~\cite{antoniadis_epta_2023}, PPTA~\cite{reardon_ppta_2023}, CPTA~\cite{xu_cpta_2023}, and the emerging IPTA~\cite{perera_ipta_2019} — demands analysis software that is physically complete, computationally efficient, and extensible to novel signal models from beyond-standard-model physics.
The existing tool PTArcade~\cite{mitridate_ptarcade_2023a} provides a user-friendly wrapper around ENTERPRISE~\cite{ellis_enterprise_2020} for MCMC-based new-physics searches.
However, no general-purpose forward simulation library exists that is purpose-built for generating large-scale simulation campaigns from diverse cosmological SGWB source models, which are the prerequisite for Simulation-Based Inference workflows~\cite{cranmer_frontier_2020}.

This proposal develops \textbf{PTAforge}: an open-source, modular, high-performance Python/C++ simulation framework for PTA SGWB signals.
PTAforge separates the physics layer (signal model: cosmic strings, phase transitions, primordial GWs, etc.), the detector layer (PTA noise, pulsar timing residuals, antenna response), and the data product layer (cross-power spectral density, Hellings-Downs correlations~\cite{hellings_upper_1983}, time-domain residuals) into independently testable, composable modules.
The framework is designed to serve as the simulation backend for SBI workflows (compatible with \texttt{saqqara}~\cite{alvey_simulationbased_2024} and \texttt{swyft}~\cite{miller_swyft_2022}) and as a standalone signal injection tool for traditional ENTERPRISE~\cite{ellis_enterprise_2020}/PTArcade~\cite{mitridate_ptarcade_2023a} analyses.

The contribution here is explicitly a software contribution: the physics is well-established, but the infrastructure to run it at scale, with reproducible outputs, modular signal composition, and GPU-accelerated batch generation, does not exist.
This addresses a real gap in the PTA software ecosystem and is publishable as a methods/software paper independent of any specific physical result.


\section{Objectives}\label{sec:3-objectives}
\begin{itemize}
    \item Design and implement the PTAforge architecture: signal module interface, detector module interface, and data product module, with clean separation of concerns and a plugin system for new signal models;
    \item Implement a library of built-in signal models: power-law SGWB~\cite{renzini_stochastic_2022}, cosmic string (VOS/BOS/LRS)~\cite{martins_extending_2002,blanco-pillado_number_2014,lorenz_cosmic_2010}, first-order phase transition~\cite{caprini_detecting_2020}, and primordial inflation as baseline plugins;
    \item Achieve GPU-accelerated batch simulation via JAX~\cite{bradbury_jax_2018} or PyTorch backends, targeting $\mathcal{O}(10^4)$ simulations per minute on a single GPU node, enabling practical SBI~\cite{cranmer_frontier_2020} simulation campaigns;
    \item Validate PTAforge outputs quantitatively against PTArcade~\cite{mitridate_ptarcade_2023a}/ENTERPRISE~\cite{ellis_enterprise_2020} results on shared benchmarks using the NANOGrav 15yr dataset~\cite{afzal_nanograv_2023,johnson_nanograv_2024};
    \item Release as a fully documented open-source package with continuous integration, contributing signal modules back to the \texttt{saqqara}~\cite{alvey_simulationbased_2024} ecosystem.
\end{itemize}

\section{Methods}\label{sec:3-methods}

\begin{itemize}
    \item The software architecture is designed following the modular pattern established by \texttt{saqqara}~\cite{alvey_simulationbased_2024}: a \texttt{SignalModel} abstract base class, a \texttt{NoiseModel} abstract base class, and a \texttt{Simulator} class that composes them and produces batched outputs.
    \item All computationally intensive inner loops (loop number density integrals, power spectrum evaluations, PTA response functions) are implemented in JAX~\cite{bradbury_jax_2018} for JIT compilation and automatic batching via \texttt{vmap}.
    \item A Python plugin registry allows user-defined signal models to be registered and used seamlessly within the SBI~\cite{cranmer_frontier_2020} or MCMC workflow without modifying core library code.
    \item Unit tests, integration tests, and numerical benchmarks are implemented via pytest and run in CI. Outputs for the power-law SGWB case are compared against PTArcade's~\cite{mitridate_ptarcade_2023a} known posteriors to validate correctness.
    \item Documentation is built with Sphinx and includes worked examples for both SBI (\texttt{swyft}~\cite{miller_swyft_2022} integration) and MCMC (PTArcade~\cite{mitridate_ptarcade_2023a}/ENTERPRISE~\cite{ellis_enterprise_2020} integration) workflows.
    \item The cluster is used for performance benchmarking of the GPU batch generation and for running the validation simulation campaigns.
\end{itemize}